\documentclass[10pt,a4paper]{amsart}
\usepackage[margin=19mm]{geometry}
\usepackage{amsmath,amssymb,mathtools}
\usepackage[scr=rsfs,cal=euler]{mathalfa}
\usepackage{xfrac}
\usepackage[unicode,hidelinks,bookmarks=false]{hyperref}
\newcommand{\ie}{i.e.\ }
\newcommand{\cf}{cf.\ }
\newcommand{\resp}{respectively\ }
\newcommand{\Subsection}{\subsection}
\providecommand{\nobreakdash}{\nobreak\hbox{-}}
\setlength{\emergencystretch}{2em}
\multlinegap0pt
\usepackage{bm}
\usepackage{bbm}
\usepackage{dsfont}
\usepackage{xspace}
\usepackage{cancel}
\usepackage{mathtools}
\usepackage{xcolor}
\usepackage{amsthm}
\makeatletter
\@ifundefined{theorem}{\newtheorem{theorem}{Theorem}}{}
\@ifundefined{lemma}{\newtheorem{lemma}[theorem]{Lemma}}{}
\@ifundefined{proposition}{\newtheorem{proposition}[theorem]{Proposition}}{}
\makeatother
\DeclareMathOperator{\dist}{{\rm dist}}
\def\D{\mathcal{D}}
\def\N{\mathbb N}
\def\R{\mathbb R}
\def\Z{\mathbb Z}
\def\cP{\mathcal P}
\title[Projection theorems with countably many exceptions and appli]{Projection theorems with countably many exceptions and applications to the exact overlaps conjecture}
\author{Meng Wu}
\date{Source: arXiv:2503.21923v2 (CC BY 4.0)}
\begin{document}
\maketitle

\noindent\emph{Source and licence.} Projection theorems with countably many exceptions and applications to the exact overlaps conjecture. Version arXiv:2503.21923v2 (CC BY 4.0), distributed under the Creative Commons Attribution 4.0 International licence, as recorded in the arXiv licence metadata for this exact version and shown on the abstract page https://arxiv.org/abs/2503.21923v2. Full rights notice: licenses/arxiv-2503.21923-CC-BY-4.0.txt.\par\smallskip

\noindent\textit{Editorial scope.} Counted unit: the Theorem whose source label is \texttt{thm:proj-cp-distribution:1} in the article (source0.tex lines 1804--1809, source label \texttt{thm:proj-cp-distribution:1}), delivered together with the article's own complete original proof of it (source0.tex lines 1820--1986, one \texttt{proof} environment of 167 lines). No mathematics is written, repaired or re-derived: statement and proof are the article's own text line for line. The proof is detached from the statement in the source: the article states the result at lines 1804--1809 and prints its proof at lines 1820--1986, and the proof environment's own header names the statement, so the association is the article's own. Required context retained uncounted: prop:dim-projection-cp-distribution:1 (lines 1340--1351); prop:slice-cp-distribution:1 (lines 1397--1399); thm:entropy-increase-1 (lines 1438--1447); lemma:entropy of restricted meas-small delete (lines 1811--1817). Every definition, display and label of those lines is retained unchanged. Omitted from this package and not redistributed: 1--1339, 1352--1396, 1400--1437, 1448--1803, 1810--1810, 1818--1819, 1987--2442. Nothing inside a retained excerpt is omitted or altered: every retained body is delivered exactly as the article's lines are written, so no omission receipt of any kind accompanies this package. Citations: the retained excerpts carry no citation command at all, so no bibliography entry of the article is transcribed and no reference is redistributed. Reference adaptation: none was needed and none was made; every reference inside the delivered unit resolves inside it, so no unresolved reference or citation is produced. Printed slips kept literal: no printed word, symbol, hypothesis or number of the counted statement or of its proof was altered. Layout and class: the article's own class is \texttt{amsart} with packages including CJKutf8, version, amssymb, amsfonts, graphicx, tikz, color, epstopdf, babel, float, cite, pgf, geometry, hyperref; the wrapper is the lane's pinned \texttt{amsart} editorial wrapper at 10pt on A4, which restates the theorem-like environments the retained text needs and adds the article's own macro definitions that the retained text needs (\texttt{\textbackslash D}, \texttt{\textbackslash N}, \texttt{\textbackslash R}, \texttt{\textbackslash Z}, \texttt{\textbackslash cP}, \texttt{\textbackslash dist}), with extra wrapper packages (bm, bbm, dsfont, xspace, cancel, mathtools, xcolor). The delivered numbering is the unit's own. Assets: no retained span contains a figure environment, a graphics-inclusion command, a PDF-page inclusion, a table or a TikZ picture; no image file is copied into this package, the package declares no asset dependency and no image file is redistributed with it. Licence: version arXiv:2503.21923v2 of this article is distributed under the Creative Commons Attribution 4.0 International licence, as recorded in the arXiv licence metadata for this exact version and shown on the abstract page https://arxiv.org/abs/2503.21923v2. A full rights notice is delivered at licenses/arxiv-2503.21923-CC-BY-4.0.txt. Characters: every retained excerpt is pure ASCII, so the delivered engine, XeLaTeX with its default Latin Modern text font, reports no missing character.
\begin{proposition}\label{prop:dim-projection-cp-distribution:1}
Let $Q$ be an ergodic CP-distribution on $\R^2$ and $\pi\in G(2,1)$.  Let $\alpha=\dim Q$ and $\beta=\dim \pi Q$. Then for $Q$-a.e.  $\mu$,  the measures $\mu$ and $\pi \mu$ are exact dimensional with $\dim \mu=\alpha$ and $\dim \pi \mu=\beta$.   In particular, for any $\epsilon>0$,  there is  $A$ such that $Q(A)>1-\epsilon$,  and for any $ \mu\in A$,  there exist $D$ and $r_0$ with $\mu(D)>1-\epsilon$ and 
\begin{eqnarray}
 \mu(B(x,r)) & = & r^{\alpha\pm \epsilon} \  \textrm{ for } x\in D,  r\le r_0,  \label{eq:prop:dim-projection-cp-distribution:1:1}\\
\pi\mu(B(\pi(x),r)) & = & r^{\beta\pm \epsilon} \  \textrm{ for } x\in D,  r\le r_0.\label{eq:prop:dim-projection-cp-distribution:1:2}
\end{eqnarray}
Moreover, for any $\epsilon>0$, there exists $l_1\in \N$ such that for all $l\ge l_1$, there is $A$ such that $Q(A)>1-\epsilon$ and for every $\mu\in A$, we have
\begin{eqnarray}
\mu\left( x: \frac{1}{l}H\left(\mu^{\D_l(x)},\D_l\right)  =  \alpha\pm \epsilon \right) &>& 1-\epsilon, \label{eq:prop:dim-projection-cp-distribution:1:3}\\
\mu\left( x: \frac{1}{l}H\left(\pi\left(\mu^{\D_l(x)}\right),\D_l\right) =  \beta\pm \epsilon \right) &>& 1-\epsilon.\label{eq:prop:dim-projection-cp-distribution:1:4}
\end{eqnarray}
\end{proposition}

\begin{proposition}\label{prop:slice-cp-distribution:1}
Let $Q$ be an ergodic  CP-distribution on $\R^2$ and $\pi\in G(2,1)$.  Suppose $\dim \pi Q = \dim Q-\kappa$ for some $\kappa>0$. Then for $Q$-a.e.  $\mu$ and $\mu$-a.e.  $x$, the conditional measure $\mu_{\pi^{-1}(\pi(x))}$ is exact dimensional and $\dim \mu_{\pi^{-1}(\pi(x))}=\kappa$.  Moreover,  for any $\epsilon>0$, there exists $A$ with $Q(A)>1-\epsilon$ and $n_0,l_0\in \N$ such that for any $n\ge n_0$, for any $\mu\in A$, there exists $D$  with $\mu(D)>1-\epsilon$ and for each $x\in D$, the measure $\mu_{\pi^{-1}(\pi(x))}$ (viewed as a measure on $\R$) is $(n,l_0,\epsilon)$-dyadic spreading.
\end{proposition}

\begin{theorem}\label{thm:entropy-increase-1}
Given $0<\gamma<1$ and $l\in \N$, there exists $\delta>0$ such that the following holds for all $n$ large enough: Let $A\subset [0,1]\cap 2^{-n}\Z$,  $\eta\in \cP([0,1]\cap 2^{-n}\Z)$.  Suppose that $\eta(A)\ge 1/2$ and for each $x\in A$,
\begin{equation}\label{eq:thm:ent-incre:1}
\left| \left\{1\le k\le n: \eta(\D_k(x))\ge 2\eta(\D_{k+l}(x))\right\} \right| \ge (1-\gamma/2)n.
\end{equation}
Let $B\subset [0,1]\cap 2^{-n}\Z$,  and for each $a\in A$,  let $B_a\subset B$.   Suppose that $|B|\le 2^{n(1-\gamma)}$ and $|B_a|\ge |B|^{1-\delta}$ for all $a\in A$. Then we have
\[
\left| \bigcup_{a\in A}(a+B_a) \right| \ge |B|^{1+\delta}.
\]
\end{theorem}

\begin{lemma}\label{lemma:entropy of restricted meas-small delete}
Let $\mu\in \cP([0,1]^d)$. Suppose that $A\subset [0,1]^d$ such that $\mu(A)>1-\epsilon$. Then for any partition $\mathcal{A}$ of $[0,1]^d$,
\[
H(\mu_{|A},\mathcal{A}) >  H(\mu,\mathcal{A})-c\sqrt{\epsilon}\log |\mathcal{A}|,
\]
where $c$ is a constant depending only on $\R^d$.
\end{lemma}

\begin{theorem}\label{thm:proj-cp-distribution:1}
Let $Q$ be an ergodic CP-distribution on $\R^2$. Then the following set is at most countable:
\begin{equation}\label{eq:thm:proj-cp-distribution:1}
\{ \pi\in S^1: \dim\pi Q < \min(1,\dim Q)\}.
\end{equation}
\end{theorem}

\begin{proof}[Proof of Theorem \ref{thm:proj-cp-distribution:1}]
Let us suppose,  by contradiction, that the exceptional set \eqref{eq:thm:proj-cp-distribution:1}  is uncountable.  It is not difficult to check that the set \eqref{eq:thm:proj-cp-distribution:1} is a Borel set.  There exists $\kappa>0$ such that the set 
\[
E=\{ \pi\in G(2,1): \dim \pi Q<\min(1,\dim Q)-\kappa \}
\]
is uncountable and there exists a non-atomic Borel probability measure $\nu$ that is supported on a compact subset of $E$.

Applying Proposition \ref{prop:slice-cp-distribution:1} to $Q$ and $\pi\in E$ and by Egorov's theorem,  we obtain that for any $\epsilon_0>0$, there exist $l_0,  n_0\in \N$,  and $E_0\subset E$ with $\nu(E_0)>1-\epsilon_0$, and for each $\pi\in E_0$, there exists $A_\pi^0$ such that $Q(A_\pi^0)>1-\epsilon_0$ and for each $\mu\in A_\pi^0$, and $n\ge n_0$  there exists $D_{\pi,\mu}^0$ with $\mu(D_{\pi,\mu}^0)>1-\epsilon_0$ and for each $x\in D_{\pi,\mu}^0$, the measure $\mu_{\pi^{-1}(\pi(x))}$ is well defined  and $(n,l_0,\epsilon_0)$-dyadic spreading.
\iffalse
{\color{red}there exists $(z,t)\in [0,1]^2\times [0,\log2]$ such that the measure $\tilde{\mu}=S^*_t(T^*_z(\mu_{\pi^{-1}(\pi(x))}))$ satisfies for all $n\ge n_0$,
\begin{equation}\label{eq:thm-proj-cp-distribution-1:1}
\tilde{\mu}\left(y:  \frac{1}{n}|\{ 1\le k\le n:  \tilde{\mu} (\D_k(y))\ge 2\tilde{\mu} (\D_{k+l}(y))\}|>1-\epsilon_0\right)>1-\epsilon_0.
\end{equation}}
\fi

Applying Proposition \ref{prop:dim-projection-cp-distribution:1} to $Q$ and $\pi\in E$ and by Egorov's theorem,  we obtain that for any $\epsilon_1>0$,  there exists $l_1\in \N$ and $E_1\subset E$ with $\nu(E_1)>1-\epsilon_1$ such that for all  $\pi\in E_1,$ and  $l\ge l_1$, there exists $A_\pi^1$ such that $Q(A_\pi^1)>1-\epsilon_1$ and for any $\mu\in A_\pi^1$, there exists $D_{\pi,\mu}^1$ with $\mu(D_{\pi,\mu}^1)>1-\epsilon_1$ and for each $x\in D_{\pi,\mu}^1$ we have (writing $\beta_\pi=\dim \pi Q$)
\begin{eqnarray}
 & & \mu(B(x,r))= r^{\alpha\pm \epsilon_1} \  \textrm{ for }  r\le 2^{-l_1},  \label{eq:thm-proj-cp-distribution-1:2}\\
& & \pi\mu(B(\pi(x),r)) = r^{\beta_\pi\pm \epsilon_1} \  \textrm{ for } r\le 2^{-l_1}, \label{eq:thm-proj-cp-distribution-1:3}\\
& & \frac{1}{l}H\left(\mu^{\D_l(x)},\D_l\right) = \alpha\pm \epsilon_1,\label{eq:thm-proj-cp-distribution-1:4}\\
& & \frac{1}{l}H\left(\pi\left(\mu^{\D_l(x)}\right),\D_l\right) = \beta_\pi\pm \epsilon_1. \label{eq:thm-proj-cp-distribution-1:5}
\end{eqnarray}
Note that by \eqref{eq:thm-proj-cp-distribution-1:3},  for each $\pi\in E_1$ and $\mu\in A_{\pi}^1$, we necessarily have
\begin{equation}\label{eq:thm-proj-cp-distribution-1:6}
N_{r}(\pi(D_{\pi,\mu}^1)) = r^{-\beta_\pi\pm \epsilon_1} \ \textrm{ for } r\le 2^{-l_1}.
\end{equation}
Note that we have $\nu(E_0\cap E_1)>1-\epsilon_0-\epsilon_1>0$, provided $\epsilon_0$ and $\epsilon_1$ were assumed small enough.  Since $\nu$ is non-atomic, the set $E_0\cap E_1$ is uncountable. In particular,  for any $\epsilon_2>0$,   there must exist $E_2\subset E_0\cap E_1$ with $\nu(E_2)>0$ such that 
 \begin{equation}\label{eq:thm-proj-cp-distribution-1:7}
|\dim \pi Q-\dim \pi' Q | <\frac{\epsilon_2}{2} \ \textrm{ for }\pi,\pi'\in E_2.
\end{equation}

In the following, we shall show that there is a contradiction if $\epsilon_0,\epsilon_1$ and $\epsilon_2$ were assumed  small enough in terms of $\beta$ and $l_0$. 

Let us pick some $\pi_0\in E_2$ and denote  $\beta_0=\pi_0 Q$.  Then for each $\pi\in E_2$, we have 
\[
\dim \pi Q = \beta_0\pm \frac{\epsilon_2}{2} .
\]
Let $\rho_0=2^{-l_1-1}$.   
 Since $E_2$ is infinite (uncountable), there exist $\pi_1,\pi_2\in E_2$ such that 
\begin{equation}\label{eq:thm-proj-cp-distribution-1:8}
0<\dist(\pi_1,  \pi_2)\le \rho_0. 
\end{equation}
Then  we have 
\[
l:=\lfloor -\log \dist(\pi_1,\pi_2) \rfloor \ge \lfloor -\log \rho_0 \rfloor \ge l_1+1.
\]
Recall that the logarithm $\log$ is in base 2. Note that we have $\dist (\pi_1,\pi_2)=2^{-l\pm 1}$.  We have seen that (consequence of Proposition \ref{prop:dim-projection-cp-distribution:1}),  for each $\pi\in \{\pi_1,\pi_2\}$, there exists $A_{\pi}^1$ with $Q(A_\pi^1)>1-\epsilon_1$ and for each $\mu\in A_\pi^1$, there exists $D_{\pi,\mu}^1$ with $\mu(D_{\pi,\mu}^1)>1-\epsilon_1$ and for each $x\in D_{\pi,\mu}^1$ we have  \eqref{eq:thm-proj-cp-distribution-1:2},  \eqref{eq:thm-proj-cp-distribution-1:3}, \eqref{eq:thm-proj-cp-distribution-1:4}, and \eqref{eq:thm-proj-cp-distribution-1:5}.
Let 
\[
A=A_{\pi_1}^0 \cap A_{\pi_1}^1 \cap A_{\pi_2}^0 \cap A_{\pi_2}^1. 
\]
Then we have $Q(A)\ge 1-2(\epsilon_0+\epsilon_1)$.   For $\mu\in A$, let 
\[
D_\mu=D_{\pi_1,\mu}^0 \cap D_{\pi_1,\mu}^1 \cap D_{\pi_2,\mu}^0 \cap D_{\pi_2,\mu}^1.
\]

Recall that we denoted $\beta_\pi=\dim \pi Q$,   and   $\beta_0=\dim \pi_0 Q$ for some $\pi_0\in E_2$.  By \eqref{eq:thm-proj-cp-distribution-1:6} and \eqref{eq:thm-proj-cp-distribution-1:7},  we have
\begin{equation}\label{eq:thm-proj-cp-distribution-1:9}
N_{r}(\pi(D_\mu)) \le r^{-\beta_0-\epsilon_2- \epsilon_1} \ \textrm{ for } \pi\in\{\pi_1,\pi_2\},r\le 2^{-l_1}.
\end{equation}
Note that we have  $\mu(D_\mu)>1-2(\epsilon_0+\epsilon_1)$.
By Markov's inequality,  there exists $\epsilon_3=o_{\epsilon_0,\epsilon_1}(1)$ (we may take $\epsilon_3=\sqrt{2(\epsilon_0+\epsilon_1)}$) and $\tilde{D}_\mu\subset D_\mu$ with $\mu(\tilde{D}_\mu)>1-\epsilon_3$ such that for each $x\in \tilde{D}_\mu$,   we have 
\begin{eqnarray}
 & & \mu_{\pi_1^{-1}(\pi_1(x))} \textrm{ is well defined and }   \mu_{\pi_1^{-1}(\pi_1(x))}(\tilde{D}_\mu) >1-\epsilon_3, \label{eq:thm-proj-cp-distribution-1:i}\\
& & \frac{\mu(\D_l(x)\cap \tilde{D}_\mu)}{\mu(\D_l(x))} > 1-\epsilon_3. \label{eq:thm-proj-cp-distribution-1:iii}
\end{eqnarray}

Since $\pi_1\mu(\pi_1\tilde{D}_\mu)\ge \mu(\tilde{D}_\mu)>1-\epsilon_3$ and for each $x\in \tilde{D}_\mu$ we have (recall \eqref{eq:thm-proj-cp-distribution-1:3}) 
\[
\pi_1\mu(B(\pi_1(x),  r)) \le r^{\beta_0-\epsilon_2-\epsilon_1} \  \textrm{ for } r\le 2^{-l_1}, 
\]
we must have
\begin{equation}\label{eq:thm-proj-cp-distribution-1:10}
N_{r}(\pi_1(\tilde{D}_\mu)) \ge c \cdot  r^{-\beta_0+\epsilon_2+ \epsilon_1} \ \textrm{ for } r\le 2^{-l_1}
\end{equation}
for some absolute constant $c$ depending only on $\R$.  Specific to $r=2^{-l}$,  it follows from \eqref{eq:thm-proj-cp-distribution-1:10} that we can find points $z_i\in \pi_1(\tilde{D}_\mu),  1\le i\le \lfloor\frac{1}{16}\cdot c \cdot  2^{l(\beta_0-\epsilon_2- \epsilon_1)}  \rfloor$, such that 
\begin{equation}\label{eq:thm-proj-cp-distribution-1:10.1}
\dist(z_i,z_j)\ge 2^{-l+3} \ \textrm{ for } i\neq  j.
\end{equation}
Let us now fix any $i\in \{1,\ldots,  \lfloor\frac{1}{16}\cdot c \cdot  2^{l(\beta_0-\epsilon_2- \epsilon_1)}  \rfloor\}$.  
In the following, we are going to show that 
\begin{equation}\label{eq:thm-proj-cp-distribution-1:10-1}
N_{2^{-2l}}\left( \pi_2\left( \pi_1^{-1}(B(z_i,2^{-l+1})) \right) \right) \ge 2^{l(\beta_0+\delta/2)}
\end{equation}
for some $\delta>0$ depending only on $l_0$ and $\beta$.
Note that since  $\dist(\pi_1,\pi_2)= 2^{-l\pm 1}$,  in view of \eqref{eq:thm-proj-cp-distribution-1:10.1}, we infer that  whenever $i\neq j$, we have 
\[
 \pi_2\left( \pi_1^{-1}(B(z_i,2^{-l+1})) \right) \cap  \pi_2\left( \pi_1^{-1}(B(z_j,2^{-l+1})) \right) = \emptyset.
\]
Hence, once \eqref{eq:thm-proj-cp-distribution-1:10-1} is proved,  we then can deduce that 
\[
N_{2^{-2l}}\left( \pi_2(\tilde{D}_\mu)\right) \ge \sum_{i=1}^{\lfloor\frac{1}{16}\cdot c \cdot  2^{l(\beta_0-\epsilon_2- \epsilon_1)}  \rfloor} 2^{l(\beta_0+\delta/2)} \ge 2^{2l(\beta_0+\delta/4)} >  2^{2l(\beta_0+2(\epsilon_2+\epsilon_1))},
\]
provided $\epsilon_2 $ and $\epsilon_1$ were assumed small enough.  This is a contradiction to \eqref{eq:thm-proj-cp-distribution-1:9}, thus concluding the proof of Theorem \ref{thm:proj-cp-distribution:1}.

Let us proceed to the proof of \eqref{eq:thm-proj-cp-distribution-1:10-1}. Let 
\[L'=\pi_1^{-1}(z_i)\cap \tilde{D}_\mu\ \textrm{ and } \ A=\{a\in \D_l(\R^2): a\cap L'\neq \emptyset\}.\]
For each $a\in A$, let us fix $x_a\in a\cap L'$.   For each $a\in A$, let $B_a'=a\cap \tilde{D}_\mu$.  Finally, let 
\[B_a=\pi_1 (B'_a) \ \textrm{ and } \  B=\pi_1(\cup_{a\in A}B'_a)=\cup_{a\in A}B_a.\]
Observe that the set $B$ has diameter bounded by $4\cdot 2^{-l}$.  By  \eqref{eq:thm-proj-cp-distribution-1:3},  for each $y\in B$ we have
\[
\pi_1\mu(B(y,  r)) = r^{\beta_0\pm\epsilon_2\pm\epsilon_1} \  \textrm{ for } r\le 2^{-l_1}.
\]
Hence we must have 
\begin{equation}\label{eq:thm-proj-cp-distribution-1:11}
N_{2^{-2l}}(B) \le  c \cdot  2^{l(\beta_0+2\epsilon_2+2 \epsilon_1)},
\end{equation}
for some absolute constant $c$ depending only on $4$ and $\R$.  

Recall that for each $x\in \tilde{D}_\mu$,  the property \eqref{eq:thm-proj-cp-distribution-1:5} is satisfied.  Combining this with  \eqref{eq:thm-proj-cp-distribution-1:iii} and Lemma \ref{lemma:entropy of restricted meas-small delete},  we obtain that the following holds for each $a\in A$:
\[
\frac{1}{l}H\left(\pi_1\left(\left(\mu_{|\tilde{D}_\mu}\right)^{\D_l(x_a)}\right),  \D_l\right)\ge \beta_0-\epsilon_2- \epsilon_1-o_{\epsilon_3}(1). 
\]
Since the measure $\pi_1\left(\left(\mu_{|\tilde{D}_\mu}\right)^{\D_l(x_a)}\right)$ is supported on $\pi_1(B_a')=B_a$, we have
\begin{equation}\label{eq:thm-proj-cp-distribution-1:12}
N_{2^{-2l}}(B_a) \ge   c \cdot  2^{l(\beta_0-\epsilon_2- \epsilon_1-o_{\epsilon_3}(1))}
\end{equation}
for some absolute constant $c>0$ depending only on $\R$.

Now, let us consider the following set 
\[
\pi_2\left( \bigcup_{a\in A}B_a' \right).
\]
Since $\dist(\pi_1,\pi_2)\le 2^{-l+1}$ and the diameter of $B_a'$ is bounded by $2^{-l+1}$,  we have
\begin{equation}\label{eq:thm-proj-cp-distribution-1:13}
d_{\rm H}(\pi_2(B_a'),  \pi_1(B_a')+ \pi_2(x_a) -\pi_1(x_a)) \le  c\cdot 2^{-2l},
\end{equation}
where $c>0$ is some absolute constant only depending on $\R^2$.  Here $d_{\rm H}$ stands for the Hausdorff distance between subsets of $\R^1$. 
Recall that for each $a\in A$ we have
\begin{equation}\label{eq:thm-proj-cp-distribution-1:14}
\pi_1(x_a)=z_i.
\end{equation}
Hence by \eqref{eq:thm-proj-cp-distribution-1:13} we have
{\color{black}
\[
d_{\rm H}\left(\pi_2\left(\bigcup_{a\in A}B_a'\right),  \bigcup_{a\in A}\left(\pi_1(B_a')+ \pi_2(x_a) -z_i\right)\right)\le  c\cdot 2^{-2l}
\]}
Recall that we denoted $B_a=\pi_1(B_a')$.  Thus, to estimate the $2^{-2l}$-cover number of $\pi_2\left( \bigcup_{a\in A}B_a' \right)$, we only need to estimate the following
\[
N_{2^{-2l}}\left( \bigcup_{a\in A}(B_a+ \pi_2(x_a) -z_i) \right).
\]
For this, we shall use Theorem \ref{thm:entropy-increase-1}.   
First,  let us recall that we have $\mu_{\pi_1^{-1}(z_i)} (L')>1-\epsilon_3$ (recall  \eqref{eq:thm-proj-cp-distribution-1:i}) and the 
measure $\mu_{\pi_1^{-1}(z_i)}$ is $(n,l_0,\epsilon_0)$-dyadic spreading.
The set $A'=\{x_a:a\in A\}\subset L'$ and $d_{\rm H}(A',L')\le c\cdot 2^{-l}$.
The set $B$ satisfies \eqref{eq:thm-proj-cp-distribution-1:11}.  Since we initially assume $\epsilon_0,\epsilon_1,\epsilon_2$ are small enough, so that we have $\beta_0+2(\epsilon_2+ \epsilon_1) \le \beta_0+(1-\beta_0)/2<1$ and  $1-\epsilon_0>\frac{1}{2}$ and $1-2\epsilon_0>\beta_0+2(\epsilon_2+ \epsilon_1) $.  Thus we may effectively apply Theorem \ref{thm:entropy-increase-1} to conclude that  there exists $\delta>0$, only depends on $l_0$ and $\beta$ such that 
\[
N_{2^{-2l}}\left( \bigcup_{a\in A}(B_a+ \pi_2(x_a) -z_i) \right)\ge \left(N_{2^{-2l}}(B)\right)^{1+\delta}\ge 2^{l(\beta_0+\delta/2)}
\]
provided $\epsilon_1, \epsilon_2,  \epsilon_3$ are assumed small enough.
Thus we have shown that 
\begin{equation}\label{eq:thm-proj-cp-distribution-1:19}
N_{2^{-2l}}\left( \pi_2\left( \bigcup_{a\in A}B_a' \right) \right) \ge c\cdot 2^{l(\beta_0+\delta/2)}
\end{equation}
for some absolute constant $c>0$.  Note that for each $a\in A$,  we have $\pi_1(B_a')\subset B(z_i,2^{-l+1})$.  Hence in particular \eqref{eq:thm-proj-cp-distribution-1:19} implies that 
\begin{equation*}\label{eq:thm-proj-cp-distribution-1:20}
N_{2^{-2l}}\left( \pi_2\left( \pi_1^{-1}(B(z_i,2^{-l+1})) \right) \right) \ge c\cdot 2^{l(\beta_0+\delta/2)}.
\end{equation*}
This is what we aimed to prove.
\iffalse
{\color{red}This holds for all $z_i,  i\in \{1,\ldots,  \lfloor\frac{1}{4}\cdot c \cdot  2^{l(\beta_0-\tau- \epsilon_1)}   \rfloor\}$. Since $\dist(\pi_1,\pi_2)\approx 2^{-l}$, we deduce that for each $i\in \{1,\ldots,  \lfloor\frac{1}{4}\cdot c \cdot  2^{l(\beta_0-\tau- \epsilon_1)}   \rfloor\}$, the number of $j\in \{1,\ldots,  \lfloor\frac{1}{4}\cdot c \cdot  2^{l(\beta_0-\tau- \epsilon_1)}  \rfloor\}$ such that $\pi_2\left( \pi_1^{-1}(B(z_i,2^{-l+1})) \right)$ intersects $\pi_2\left( \pi_1^{-1}(B(z_j,2^{-l+1})) \right)$ is at most $8$.  Thus we obtain
\[
N_{2^{-2l}}\left( \pi_2(\tilde{D}_\mu)\right)\ge \frac{1}{8}\cdot c\cdot 2^{l(\beta_0-\tau- \epsilon_1+\beta_0+\delta/2)} >2^{2l(\beta_0+\delta/4)}
\]
provided $\tau,\epsilon_1$ are assumed small enough.  This is a contradiction to \eqref{eq:thm-proj-cp-distribution-1:9}.}
\fi
\end{proof}

\end{document}
